\section{How sparse can a one-sided Frobenius test be?}
\label{sec:transfer}
The residue argument in
\texttt{New.LFunction\_ne\_zero\_of\_primeRemainder\_le} suggests a
finite extremal question. Once a quotient of a Galois group has been
controlled, can a bound for a single nonidentity splitting type see
every remaining irreducible character with the required sign? We give
a sharp answer for tests designed by the signs of individual character
coefficients. This qualification matters because additional arithmetic
relations among zero multiplicities can allow other tests.

\subsection*{An exact support and minimax theorem}
Let $G$ be a finite group of order $g$, and let $N\ne1$ be a normal
subgroup of order $n$. Call a complex irreducible character
$\chi$ \emph{new} in this section if its representation is nontrivial
on $N$. All other irreducibles come from $G/N$. Let $f$ be a real
class function on $G$ with $f(u^{-1})=f(u)$, $0\le f\le 1$, and $f(1)=0$.
Write $f=\sum_\chi a_\chi\chi$, with the usual normalized character
inner product as coefficient. The $a_\chi$ are real by inversion
symmetry. Suppose $a_\chi\le0$ for every new $\chi$.
Define the worst response to one unit of regular-character multiplicity by
\[
  t(f)=\min_{\chi\ \mathrm{new}}\frac{-a_\chi}{\chi(1)}.
\]

\noindent\textbf{Theorem.}
The zero set of $f$ on $N$ is exactly
$N\cap\bigcap_{\chi\ \mathrm{new},\ a_\chi<0}\ker\chi$.
If every new coefficient is strictly negative, the support contains
$N\setminus\{1\}$. The indicator
$f_*=\mathbf 1_N-\mathbf 1_{\{1\}}$ attains this minimal support.
Moreover,
\[
                      \max_f t(f)=1/g.
\]
The maximizers are precisely $f_*+h$ for real inversion-symmetric
class functions $h$ inflated from $G/N$, with $h(N)=0$ on the identity
coset and with values chosen so that $0\le f_*+h\le1$.
Thus $f_*$ is the unique maximizer with minimal support.

\smallskip\noindent\emph{Proof.}
For $u\in N$ the characters inflated from $G/N$ take their full degree
there. Subtracting $f(1)=0$ and taking real parts gives
\[
  f(u)=\sum_{\chi\ \mathrm{new}}(-a_\chi)
                 \bigl(\chi(1)-\operatorname{Re}\chi(u)\bigr).
\]
Each bracket is nonnegative. Equality in a bracket holds exactly when
the unitary matrix representing $u$ is the identity, by summing the real
parts of its eigenvalues. This proves the zero-set statement.
The regular representation is faithful, so all new irreducibles
together separate the identity from every other element of $N$.
The mean of a new character over $N$ is zero, since averaging its
representation over $N$ projects to its $N$-invariants, a proper
$G$-submodule and hence zero. The squared degrees of the new
irreducibles sum to $g-g/n$. We obtain
\[
 t(f)(g-g/n)
 \ \le\ \sum_{\chi\ \mathrm{new}} (-a_\chi)\chi(1)
 \ =\ \frac1n\sum_{u\in N} f(u)\ \le\ 1-1/n.
\]
As $a_\chi(f_*)=-\chi(1)/g$ for each new $\chi$, both estimates are
sharp. Equality in the first requires every new coefficient to have
this value. The difference $f-f_*$ then inflates from the quotient,
and $f(1)=f_*(1)=0$ forces the asserted vanishing on $N$. The converse
follows from the same coefficient calculation.\hfill$\square$

A consequence in the opposite sign convention is worth noting. If
$f(1)=0$, $f\ge0$, and all new coefficients are nonnegative, the same
identity, with its signs reversed, forces all of them to vanish
(average over $N$). No strictly positive response of that orientation
is possible under these conditions.

\subsection*{Exactly when one nonidentity rational Frobenius type suffices}
Now take $N=D=[G,G]\ne1$. The known quotient is the abelianization.
A rational conjugacy class here means the generators of all conjugates
of one cyclic subgroup. Then a test supported on one nonidentity
rational conjugacy class has strictly negative coefficients on all
nonlinear irreducibles if and only if $D\setminus\{1\}$ is one rational
conjugacy class. The parallel statement holds for ordinary conjugacy
classes. Indeed, a class containing any member of $D$ is contained in
$D$ in either sense, and the support theorem forces it to contain
$D\setminus\{1\}$. The converse uses $f_*$.

There is a structural characterization. In the rational case $D$ must be
the additive group of $\mathbb F_{p^d}$ for some prime $p$, and, under
a suitable identification, the conjugation image $H$ of $G$ on $D$
must consist of field multiplications and satisfy
$H\mathbb F_p^\times=\mathbb F_{p^d}^{\times}$. In the ordinary case
the condition is $H=\mathbb F_{p^d}^{\times}$.
To prove necessity, all nonidentity elements of $D$ have the same
order, necessarily a prime $p$ by taking powers. Thus $D$ is a
$p$-group, and its nontrivial center is characteristic, invariant
under conjugation and taking unit powers. Transitivity forces
$D=Z(D)$, hence $D$ is an elementary abelian vector space $V$.
The conjugation image $H$ is abelian because $G/D$ is abelian.
The group generated by $H$ and the prime-field scalar maps acts
transitively on $V\setminus0$, so $V$ is irreducible as an
$\mathbb F_p[H]$-module. The commutative algebra generated by $H$
in $\operatorname{End}_{\mathbb F_p}(V)$ is a finite field: every
nonzero element has invariant kernel and is invertible, and its inverse
belongs to the algebra by finiteness. Irreducibility makes $V$ one
dimensional over this field. The transitivity statements now give the
claimed scalar conditions, which also prove sufficiency.

For $G=U_3(\mathbb F_p)$, the upper unitriangular $3$ by $3$ matrices,
$D$ is its center, the additive group supported in the $(1,3)$ entry.
Indeed $[(a,b,c),(a',b',c')]$ has central coordinate $ab'-a'b$ in
coordinates $(a,b,c)=(u_{12},u_{23},u_{13})$. The nonzero center is
one rational class, consisting of $p-1$ ordinary classes. Its density
is $(p-1)/p^3$. For $p>2$, observing one ordinary class cannot meet
the uniform sign and coverage requirements of the theorem.

\subsection*{The arithmetic interpretation needs only Dedekind zeta functions}
Let $L/\mathbb Q$ be finite Galois with group $G$, let $F=L^N$, and set
$a=[F:\mathbb Q]=g/n$. Assume $\zeta_F$ has no zeros in
$\operatorname{Re}s>1/2$. For $N=D$ this is supplied by the
Dirichlet theorem of the specified repository and the abelian
factorization of $\zeta_F$. The event $f_*(\operatorname{Frob}_p)=1$
means that the unramified prime $p$ splits completely in $F$ and fails
to split completely in $L$. Its density is $\delta=1/a-1/g$. Put
\[
  R(x)=\sum_{\substack{p\le x,\ p\ \mathrm{unramified}\\
                          f_*(\operatorname{Frob}_p)=1}}
       (\log p)(p^{-1}-x^{-1})-\delta\log x-c.
\]
Here $c$ is the real constant specified below. Then the following
conditions are equivalent:
\begin{enumerate}
\item Every nontrivial zero of $\zeta_L$ has real part $1/2$.
\item For each $0<\epsilon<1/2$, there is $C_\epsilon$ such that
$R(x)\ge-C_\epsilon x^{-1/2+\epsilon}$ for all sufficiently large $x$.
\end{enumerate}

\noindent\emph{Proof.}
Let $B(s)$ be the series over the same primes with summands
$(\log p)p^{-s}$. In $\operatorname{Re}s>1$ the unramified
prime-power coefficients of $-\zeta_M'/\zeta_M$ for a Galois field $M$
are $[M:\mathbb Q](\log p)\mathbf 1_{\operatorname{Frob}_p^k|_M=1}$
at $p^k$. Thus
\[
 B(s)=-\frac{1}{a}\frac{\zeta'_F}{\zeta_F}(s)
             +\frac{1}{g}\frac{\zeta'_L}{\zeta_L}(s)+H(s),
\]
where $H$ is holomorphic for $\operatorname{Re}s>1/2$. It removes
higher prime powers by absolute convergence and removes the finitely
many ramified Euler factors. If $b$ is the constant coefficient of
$B(s)-\delta/(s-1)$ at $1$, set $c=b-\delta$. Termwise integration for
$\operatorname{Re}s>1$ gives
\[
 \int_1^\infty R(x)x^{-s}\,dx=
 \frac{B(s)}{s(s-1)}-\frac{\delta}{(s-1)^2}-\frac{c}{s-1}.
\]
Our choice of $c$ removes the pole at $1$. In $1/2<\operatorname{Re}s<1$
the only remaining poles come from zeros $\rho$ of $\zeta_L$, with
residue $\operatorname{ord}_\rho\zeta_L/(g\rho(\rho-1))$.
In particular each real pole has negative residue.

Assume the lower bound and set $\beta=1/2+\epsilon$. The function
$V(x)=R(x)+C_\epsilon x^{\beta-1}$ is eventually nonnegative.
Changing it on a bounded interval adds an entire function to its
transform. If the abscissa of convergence of this nonnegative transform
were strictly larger than $\beta$, Landau's theorem would force a
singularity at that real abscissa. There is no such singularity at or
above $1$. Below $1$ a real pole would have negative residue, impossible
for a nonnegative integral approached from the right. Consequently
the integral is holomorphic in $\operatorname{Re}s>\beta$, and the
displayed residue forbids any $\zeta_L$ zero there. Letting $\epsilon$
decrease to zero and using the functional equation gives (1).

Conversely, the prime ideal theorem under RH for each of $F,L$ gives
$\psi_M(x)=x+O_{M,\eta}(x^{1/2+\eta})$ for every $\eta>0$.
Removing the finitely many ramified primes and higher powers
costs $O_L(\sqrt{x}\log^2(2x))$. Therefore the sum $P(x)$ of $\log p$
over the selected primes is $\delta x+O_{L,\eta}(x^{1/2+\eta})$.
Partial summation identifies the weighted sum with
$\int_1^xP(t)t^{-2}\,dt$. Its remainder after subtracting
$\delta\log x$ tends to a constant with error
$O_{L,\eta}(x^{-1/2+\eta})$ for $0<\eta<1/2$. The Mellin formula
at $1$ identifies this constant with $c$ (multiply the transform
by $s-1$ and let $s\downarrow1$).\hfill$\square$

\paragraph{An explicit quartic-nonresidue test.}
Take $L=\mathbb Q(i,\sqrt[4]{2})$. The splitting field has degree $8$:
$X^4-2$ is Eisenstein, and the real field $\mathbb Q(\sqrt[4]{2})$
does not contain $i$. A generator that multiplies $\sqrt[4]{2}$ by
$i$ while fixing $i$, and complex conjugation, give the dihedral group
of order $8$. The commutator has order $2$ and fixes
$F=\mathbb Q(i,\sqrt2)=\mathbb Q(\zeta_8)$.
The event for an odd prime is exactly
\[
 p\equiv1\pmod8,\qquad 2^{(p-1)/4}\equiv-1\pmod p.
\]
Indeed the first condition gives splitting in $F$. For these primes
$2$ is a square modulo $p$ and all fourth roots of unity are available,
so that failure of complete splitting in $L$ is precisely the quartic
nonresidue condition. The first examples are $17,41,97,137,193$.
The density and the optimal worst coefficient are both $1/8$ here.
An estimate from below for the displayed smoothed count, with main
term $(\log x)/8$ and its own constant, would settle RH for this
nonabelian splitting field. Establishing that lower estimate is the
unresolved analytic step. Neither the support theorem nor a finite
prime count supplies it.

\paragraph{Related literature.}
Suzuki studies one-sided criteria for RH and for Dirichlet GRH in
\emph{On variants of Chebyshev's conjecture},
\emph{Ramanujan Journal} 68 (2025), article 95
(\url{https://arxiv.org/abs/2411.07436}).
Bella\"iche studies character coefficients in effective Chebotarev
bounds (\url{https://arxiv.org/abs/1308.1022}).
Lewis, Morotti, Pacifici, Sanus and Tong-Viet treat ordinary conjugacy
classes equal to a difference of normal subgroups
(\url{https://arxiv.org/abs/2603.25407}).
The finite extremum above addresses a specified normalized coefficient
and a specified sign cone. Its proof is elementary character theory.
We make no claim of historical priority for the criterion or its group
theory.
